\begingroup\sloppy
\providecommand{\F}{\mathbb{F}_p}
\providecommand{\enc}{\mathrm{Enc}}
\providecommand{\T}{\top}
\providecommand{\sC}{\mathsf{C}}
\providecommand{\sK}{\mathsf{K}}
\providecommand{\Kpre}{\mathsf{K}_{\mathrm{pre}}}
\providecommand{\para}[1]{\par\noindent\textbf{#1}\quad}
\providecommand{\todo}[1]{\textbf{[TODO: #1]}}
\providecommand{\Fext}{\mathbb{K}}
\providecommand{\Vlo}{\mathrm{lo}}
\providecommand{\Vshift}{\mathrm{shift}}
\providecommand{\Vrq}{\mathrm{rq}}
\providecommand{\Vtab}{\mathcal{T}}
\providecommand{\Vcut}{\mathcal{C}}
\providecommand{\Vclr}{\mathcal{O}}
\providecommand{\Nleaf}{N_{\mathrm{leaf}}}
\providecommand{\Veq}{\mathrm{eq}}
\providecommand{\toprule}{\hline}
\providecommand{\midrule}{\hline}
\providecommand{\bottomrule}{\hline}

\section{Proofs of the Optimisations}
\label{app:opt}

\subsection{Conventions for the Proofs of the Optimisations}
\label{app:the-conv}

This part of the appendix proves that the optimisations of our implementation, on which
Sections~\ref{sec:defence} to~\ref{sec:gen} rely, keep the soundness theorems: last-block pruning,
generation as one prefill with the token rule and the integer sampler, the Fiat--Shamir transcript
with tree absorption of large messages and its omission in mode $\Kpre$, the compact encoding of the
proof, and mode $\Kpre$ with V1 and with secret vectors reused over many queries. It also proves the
acceptance bound of the path protocol for edits of $k$ neurons. Each proof names the functions of
the library \texttt{code/src/pvi} that it models.

\para{Graphs.} A graph $\mathcal G$ is the ordered list of weight operations~\eqref{eq:weightop} and
operations without weights of Section~\ref{sec:intmodel}, and $M_{\mathcal G}(x)$ is its output on
the query $x$: the claims of its last weight operation, for a decoder the logits. Every statement
about values is a statement about the integer semantics of Appendix~\ref{app:semantics}. That the
optimised kernels compute these semantics (the int32 and int64 attention paths, the native CPU
kernels, the fused GPU kernels) is a property of the implementation that our tests check and that
we do not prove.

\para{Verdicts and query counts.} After a query the prover learns its \emph{verdict}: acceptance,
or the label of the first check that failed (\texttt{range\_or\_shape}, \texttt{token}, the lookup
checks, \texttt{freivalds}, the column checks, and for V1 \texttt{cut\_exception},
\texttt{cut\_multiplicities}, \texttt{cut\_gkr}, \texttt{kpre\_y}, \texttt{cut\_final},
\texttt{cut\_logup}). We write $Q_V$ for the number of verdicts the prover obtains for one key of
mode $\Kpre$, and $Q_H$, or $Q$ where no confusion arises, for the number of its random-oracle
queries. Unless stated otherwise, nothing else that depends on the verifier's secrets, such as its
running time, reaches the prover (Remark~\ref{rem:stateful}).

\para{Hash functions and challenges.} We model SHAKE-256~\cite{fips202} as a random oracle $H$ with
outputs of any length, and write $H_{256}$ for its first 32 bytes. An interactive challenge is
expanded with $H$ from a fresh 256-bit key that the verifier draws from the operating system
(\texttt{Challenger.\_key}). The expansion is exact: field elements come from 32-bit words below
$2p$ reduced modulo $p$, which are uniform on $\F$ since $p$ divides $2p$
(\texttt{Challenger.folding}, \texttt{ext}), and column sets from 64-bit words below the largest
multiple of $n$ that does not exceed $2^{64}$, reduced modulo $n$, with repeats discarded; the
first $t$ distinct values of an independent uniform sequence form a uniform set of $t$ distinct
columns (\texttt{Challenger.columns}). As in Section~\ref{sec:protocol} we treat a
challenge expanded from a key the prover has not seen as uniform. In the random-oracle model this
costs at most $Q_H\,2^{-256}$ per key: the prover's view depends on the expansion only if one of its
queries begins with the 256-bit key.

\para{Round-by-round soundness.} A public-coin protocol is round-by-round sound with per-round
errors $\varepsilon_1,\dots,\varepsilon_k$~\cite{canetti2019fs} if there is a set of \emph{doomed}
partial transcripts that contains the empty transcript of every false statement, that no prover
message leaves, whose complete members the verifier rejects, and that the $i$-th challenge leaves
with probability at most $\varepsilon_i$. We use three consequences, in the form and numbering of
the textbook of Chiesa and Yogev~\cite{chiesa2024snargs} (version of 25 March 2026):
(R1)~the soundness error is at most $\sum_i\varepsilon_i$ (Claim~31.1.3);
(R2)~with $\varepsilon_{\mathrm{rbr}}=\max_i\varepsilon_i$, a state-restoration prover with $t$
moves wins with probability at most $(t+k)\,\varepsilon_{\mathrm{rbr}}$ (Theorem~31.2.1; Holmgren
shows that the two notions are equivalent~\cite{holmgren2019rbr}); and
(R3)~the Fiat--Shamir transformation with the full transcript as the hash input inherits the
state-restoration error (Theorem~14.3.1; see also~\cite{bcs16}).

\subsection{Last-Block Pruning}
\label{app:prune}

\para{The pruned graph.} Let $\mathcal G$ be a decoder of $L$ blocks built by
\texttt{pvi.models.decoder.build\_decoder}, and $\mathcal G^{(n)}$ (\texttt{with\_lm\_positions}) its
form whose LM head reads the last $n\ge1$ positions, so that $\mathcal G=\mathcal G^{(1)}$. For a
tensor $X$ with positions $0,\dots,T-1$ (a column per position for the operand of a weight
operation, a row per position for an activation) let $S_n(X)$ keep the positions $T-n,\dots,T-1$.
Block $L$ of $\mathcal G^{(n)}$ computes at every position $h=\mathrm{norm}_1(r)$; $q$, $k$ and $v$
from $h$ by weight operations and requantisations, an optional per-head norm, and, in models with
rotary embeddings, RoPE of $q$ and $k$; the attention output $a=\mathrm{Att}(q,k,v)$, requantised; $r'=\mathrm{res}(r,W_o a)$; then
$h'=\mathrm{norm}_2(r')$, the MLP and $r''=\mathrm{res}(r',W_{\mathrm{down}}g)$. The LM head then
reads $S_n(\mathrm{norm}_f(r''))$. The pruned form $\mathcal G^{(n)}_{\mathrm{pr}}$
(\texttt{prune\_last=True}) computes $q$ from $S_n(h)$, applies RoPE to it at the positions
$T-n,\dots,T-1$ (\texttt{\_rope} with an offset), lets only these $n$ queries attend
(\texttt{integer\_attention} with $T_q=n$ query rows and the last $n$ rows of the causal mask), forms
$\mathrm{res}(S_n(r),W_o a_{\mathrm{pr}})$, and runs the rest of the block at these $n$ positions;
$k$ and $v$ are still computed at every position. The builder calibrates the pruned block on every
position (\texttt{\_Builder.last} with \texttt{calibrate\_all}), so both graphs have the same
weights, multipliers and tables (asserted in \texttt{tests/test\_pruning.py}).

\begin{definition}[position-local operations]
An operation $f$ without weights is \emph{position-local} if there are functions $f_0,f_1,\dots$
such that, for every $T$ and every $t<T$, position $t$ of $f(X^1,\dots,X^m)$ is
$f_t(X^1_t,\dots,X^m_t)$; it is \emph{position-uniform} if moreover all $f_t$ are equal. For a
weight operation, $Z_\ell=A_\ell[X_\ell;\mathbf 1]$ has column $t$ equal to $A_\ell[X_{\ell,t};1]$,
so it is position-uniform. \emph{Causal attention} is
$\mathrm{Att}(Q,K,V)_t=g\big(q_t,(k_j,v_j)_{j\le t}\big)$ for a function $g$ that does not depend on
$T$, and for an offset $o$ we write $\mathrm{Att}_o(Q',K,V)_i=g\big(q'_i,(k_j,v_j)_{j\le o+i}\big)$.
\end{definition}

\begin{lemma}[locality of the decoder]
\label{lem:prune-local}
In the integer semantics of Appendix~\ref{app:semantics}, every operation of the decoders of
\texttt{build\_decoder} other than attention and the slices $S_n$ is position-local. All are
position-uniform except RoPE and the learned position ids, whose functions $f_t$ depend on $t$ only
through the absolute position index. Attention is causal attention.
\end{lemma}

\begin{proof}
Embedding look-ups read the row of the token at position $t$; the learned position ids are $t$
itself. Requantisations, residual additions, the activation tables and SwiGLU act entrywise. Both
norms, and the per-head norm of $q$ and $k$, act on the channels of one position. RoPE rotates the
pairs of position $t$ by the angles $t\theta_j$, which depend on the absolute index $t$ only. For
attention, row $t$ forms the scores $s_j=q_t\cdot k_j$, the gaps $\max_{j'\le t}s_{j'}-s_j$ over the
keys $j\le t$ that the causal mask allows, $e_j$ from the public table and $e_j=0$ for $j>t$, then
$S=\sum_je_j$ and $P_j=\lfloor(2P_{\max}e_j+S)/(2S)\rfloor$, with the public constant $P_{\max}$
(255 in the cost benchmarks, 2{,}047 in the certified models), which is $\lfloor S/(2S)\rfloor=0$ for
$j>t$.
The output $\sum_jP_jv_j$, requantised entrywise, therefore depends only on $q_t$ and on
$(k_j,v_j)_{j\le t}$, through formulas that do not involve $T$. (The implementation chooses between
an int32 and an int64 path depending on $T$; both compute these integers.)
\end{proof}

\begin{lemma}[slicing commutes]
\label{lem:prune-slice}
For a position-local $f$ and $o=T-n$, $S_n(f(X^1,\dots,X^m))=f^{[o]}(S_n(X^1),\dots,S_n(X^m))$,
where $f^{[o]}$ applies $f_{o+i}$ at its position $i$; for a position-uniform $f$ this is
$f(S_n(X^1),\dots,S_n(X^m))$. For causal attention,
$S_n(\mathrm{Att}(Q,K,V))=\mathrm{Att}_{o}(S_n(Q),K,V)$.
\end{lemma}

\begin{proof}
Position $o+i$ of $f(X^1,\dots,X^m)$ is $f_{o+i}(X^1_{o+i},\dots,X^m_{o+i})$, and $X^l_{o+i}$ is
position $i$ of $S_n(X^l)$. Row $o+i$ of $\mathrm{Att}(Q,K,V)$ is
$g(q_{o+i},(k_j,v_j)_{j\le o+i})$, and $q_{o+i}$ is row $i$ of $S_n(Q)$.
\end{proof}

\begin{theorem}[last-block pruning]
\label{thm:prune}
Let $n\ge1$ and let $x$ be a query with $T\ge n$ positions.
(a)~$M_{\mathcal G^{(n)}_{\mathrm{pr}}}(x)=M_{\mathcal G^{(n)}}(x)$. Moreover, every tensor of
$\mathcal G^{(n)}_{\mathrm{pr}}$ equals the corresponding tensor of $\mathcal G^{(n)}$ (in blocks
$1$ to $L-1$, and $h$, $k$ and $v$ of block $L$) or its slice $S_n$ (the other tensors of block
$L$, and the final norm); in particular the honest claims of every weight operation of $\mathcal G^{(n)}_{\mathrm{pr}}$
are those of $\mathcal G^{(n)}$ or their last $n$ columns.
(b)~Every weight operation of $\mathcal G^{(n)}_{\mathrm{pr}}$ meets the setup bound of
Section~\ref{sec:intmodel} if and only if it does in $\mathcal G^{(n)}$.
(c)~Theorems~\ref{thm:sound}, \ref{thm:untrusted} and~\ref{thm:v1} hold for
$\mathcal G^{(n)}_{\mathrm{pr}}$ with its own commitment, set $\mathcal J$ and cut set. In
particular, the verifier of $\mathcal G^{(n)}_{\mathrm{pr}}$ in mode $\sC$ accepts an output
$y\neq M_{\mathcal G^{(n)}}(x)$ with probability at most the bound $\varepsilon_{\sC}$ computed for
the commitment of $\mathcal G^{(n)}_{\mathrm{pr}}$. If every weight operation is committed in a
row-layout tree of its own, the two graphs have the same $\mathcal J$, $k_j$, $n_j$ and $t_j$, and
therefore the same bound.
\end{theorem}

\begin{proof}
(a)~Blocks $1$ to $L-1$ and the embeddings are the same operations on the same input, and so is
$h=\mathrm{norm}_1(r)$ of block $L$. By Lemmas~\ref{lem:prune-local} and~\ref{lem:prune-slice},
the pruned $q$ is $\mathrm{rq}(W_q[S_n(h);\mathbf 1])=S_n(\mathrm{rq}(W_q[h;\mathbf 1]))$, its
per-head norm commutes with $S_n$, and RoPE at the positions $o+i$, $o=T-n$, gives
$S_n(\mathrm{RoPE}(q))$, which is $f^{[o]}$ of Lemma~\ref{lem:prune-slice}. With $\tilde q$,
$\tilde k$, $\tilde v$ the inputs of attention in $\mathcal G^{(n)}$, the pruned attention is
$\mathrm{Att}_o(S_n(\tilde q),\tilde k,\tilde v)=S_n(\mathrm{Att}(\tilde q,\tilde k,\tilde v))$, and
the output requantisation and $W_o$ are position-uniform, so the pruned $W_oa$ is $S_n(W_oa)$. The
residual addition gives $\mathrm{res}(S_n(r),S_n(W_oa))=S_n(r')$. Every remaining operation of the
block is position-uniform, so the pruned block ends with $S_n(r'')$. The pruned graph then applies
$\mathrm{norm}_f$ and its slice to a tensor of $n$ positions, on which $S_n$ is the identity, so the
LM head reads $\mathrm{norm}_f(S_n(r''))=S_n(\mathrm{norm}_f(r''))$, its input in
$\mathcal G^{(n)}$. The weights are equal, so are the logits. The statement on the claims is the
same chain of equalities.
(b)~The bound checked at setup (Section~\ref{sec:intmodel}; \texttt{pvi.v0.params.claim\_bound})
depends only on the weights and on the input bound of the operation, which are the same in both
graphs.
(c)~The proofs of the three theorems (Appendices~\ref{app:sound}, \ref{app:untrusted}
and~\ref{app:v1}) use only that the graph consists of weight operations~\eqref{eq:weightop},
committed as in Section~\ref{sec:commit}, and of public deterministic operations without weights
that the verifier recomputes, together with the setup bound. The pruned graph is such a graph: the
slices are operations without weights, and the graph digest of Section~\ref{sec:fs} describes them
(\texttt{IntGraph.digest} hashes every operation's name, inputs, output, note, function and
constants, and the graph's \texttt{meta}), so the statements of the two graphs differ under
Fiat--Shamir. The theorems therefore hold for $\mathcal G^{(n)}_{\mathrm{pr}}$, and by~(a) the event
$y\neq M_{\mathcal G^{(n)}_{\mathrm{pr}}}(x)$ is the event $y\neq M_{\mathcal G^{(n)}}(x)$. If every
weight operation has its own row-layout tree, both graphs commit the same matrices with the same
shapes and codeword lengths, so $\mathcal J$, $k_j$, $n_j$ and $t_j$ agree. The roots themselves
differ: each tree is domain-separated by its operation's name, and the builder numbers operations in
order, so the extra slices of the pruned block shift the names after them. A client therefore needs
the commitment of the graph it verifies, which it recomputes in the same way. Under a commitment
plan the plan may also differ: in the pruned graph $q$ of block $L$ reads $S_n(h)$ and $k$ and $v$
read $h$, so they are not stacked into one transposed matrix (\texttt{tests/test\_pruning.py},
\path{test_a_pruned_block_commits_its_q_on_its_own}), and the bound is that of the pruned
graph's plan.
\end{proof}

\noindent For V1, the operations of the pruned block that claim $n<T$ columns violate rule~(P2) of
Section~\ref{sec:v1} and are clear; Theorem~\ref{thm:v1} applies with this cut set.
The real-weight builders (\texttt{pvi.models.hf\_opt.build\_opt\_from\_hf} and
\texttt{pvi.models.hf\_llama.build\_llama\_from\_hf}) insert the same slices with \texttt{prune\_last}; their
graphs have the operations and order of \texttt{build\_decoder}'s (\texttt{matches\_benchmark\_graph}),
and their next-token logits with and without pruning are tested to be equal, so the theorem applies to
them. A fixed public prefix (\texttt{prefix\_tokens}) enters only as constants of the attention
operation.

\subsection{Generated Responses: the Token Rule and the Integer Sampler}
\label{app:gen}

\para{Responses and rules.} A prompt $x=(x_1,\dots,x_P)$ is answered by a response
$y_1,\dots,y_G$. A \emph{response rule} maps a logit vector $z\in\mathbb Z^V$ and a position index
$\pi$ to a token. The greedy rule is $R(z,\pi)=\min\{v:z_v=\max_wz_w\}$. The sampled rule
$R_\sigma$, for $\sigma=(\mathit{seed},a,k)$ with a 32-byte seed, $0<a<2^{24}$ and $k\ge0$
(\texttt{pvi.v0.token\_sampler.TokenSampler}), takes the candidates $C$, which are all $v$ if $k=0$ or
$k\ge V$ and otherwise the $k$ indices of largest $z_v$ with ties to the smaller index; the noise
$\Gamma[u_v]$ of Eq.~\eqref{eq:gumbel-table}, where $u_v$ is the $v$-th little-endian 16-bit word of
$H(\texttt{pvi/sampler/v1}\,\|\,\mathit{seed}\,\|\,\mathrm{le}_8(\pi))$ and $\mathrm{le}_8(\pi)$ is
$\pi$ in 8 bytes, little-endian (\texttt{TokenSampler.noise}); and returns $\min\{v\in C:\ az_v+\Gamma[u_v]\text{ is maximal on
}C\}$. The \emph{reference response} of $x$ under a rule $R$ is $y^\ast$ with
$y^\ast_i=R(\ell_i,P+i-2)$, where $\ell_i=M_{\mathcal G}(x\,\|\,y^\ast_1\cdots y^\ast_{i-1})$ are the
logits at the last position, whose 0-based index is $P+i-2$ (\texttt{greedy\_tokens},
\texttt{sample\_tokens}).

\para{The query.} The response is proved as one query on $s=x\,\|\,y_1\cdots y_{G-1}$, with
$T=P+G-1$ positions, on the graph $\mathcal G^{(G)}$ (pruned or not). Column $j<G$ of the claimed
LM-head output $\hat Z\in\mathbb Z^{V\times G}$ belongs to position $P-1+j$. The client forms $s$
from its own prompt and the received tokens. The token rule (\texttt{Verifier.check\_tokens})
requires $y_{j+1}=R(\hat Z_{:,j},P-1+j)$ for every $j<G-1$, and the client takes
$y_G=R(\hat Z_{:,G-1},P+G-2)$.

\begin{lemma}[causal prefixes]
\label{lem:causal}
For every query $s$ with $T$ positions, every $n\le T$ and every $j<n$, column $j$ of
$M_{\mathcal G^{(n)}}(s)$ equals $M_{\mathcal G}(s_0\cdots s_{T-n+j})$, the last-position logits
of $\mathcal G$ on the prefix of length $T-n+j+1$. The same holds for
$\mathcal G^{(n)}_{\mathrm{pr}}$.
\end{lemma}

\begin{proof}
By Lemma~\ref{lem:prune-local} every operation is position-local or causal attention, and none has
a constant that depends on the input. By induction over the operations in graph order, every tensor
at position $t$ computed on $s$ equals the same tensor at position $t$ computed on any prefix
$s_0\cdots s_{t'}$ with $t\le t'$: a position-local operation reads position $t$ of its inputs and
the index $t$, and row $t$ of causal attention reads rows $j\le t$ of its inputs. Column $j$ of the
LM head of $\mathcal G^{(n)}$ is position $T-n+j$, which is the last position of the prefix of
length $T-n+j+1$, where $\mathcal G$ applies its LM head. For the pruned graph, combine this with
Theorem~\ref{thm:prune}(a).
\end{proof}

\begin{lemma}[the sampler is exact]
\label{lem:sampler}
For every $z$ with $|z_v|<2^{29}$, every $\sigma$ and every $\pi$, all scores $az_v+\Gamma[u_v]$
are integers of magnitude below $2^{54}$, which 64-bit arithmetic computes exactly, and the
candidates and the choice among equal scores are determined. Hence $R_\sigma(z,\pi)$ is a function
of $(z,\sigma,\pi)$ that every platform computes identically, given the public table $\Gamma$.
\end{lemma}

\begin{proof}
$|az_v|<2^{24}\cdot2^{29}=2^{53}$, and $\Gamma$ is non-decreasing with end values
$\Gamma[0]=-161{,}658$ and $\Gamma[2^{16}-1]=772{,}243$ (Section~\ref{sec:ext-sampling}), so
$|\Gamma[i]|<2^{20}$ and every score lies below $2^{53}+2^{20}<2^{54}$ in magnitude. The
implementation gives excluded entries the score $-2^{63}$, below every candidate's, chooses the
candidates by a stable sort of $-z$ (ties to the smaller index) and the token by
\texttt{torch.argmax}, which returns the first maximal index. The table is fixed by its SHA-256
digest: \texttt{gumbel\_table} checks the table computed in floating point against the digest,
recomputes it in decimal arithmetic if they differ, and stops if the digest still does not match.
\end{proof}

\noindent The token rule runs before the range check of the claims in \texttt{run\_query}; since
acceptance requires the range check, $|z_v|<2^{29}$ holds for the claims of every accepted query,
and the lemma applies to them.

\begin{theorem}[generated responses]
\label{thm:gen}
Fix a prompt $x$, a rule $R$ (greedy, or $R_\sigma$) and $G\ge1$. For every prover, the verifier
accepts and outputs a response different from the reference response $y^\ast$ with probability at
most $\varepsilon$, the bound of Theorem~\ref{thm:sound} (or of Theorem~\ref{thm:v1} for V1) for the
graph $\mathcal G^{(G)}$ in the mode used. Under Fiat--Shamir the bound of Theorem~\ref{thm:fstree}
applies to the statement $(\mathcal G^{(G)},s,\sigma)$; in mode $\Kpre$, a false response among
$Q_V$ verdicts is accepted with probability at most $Q_V\varepsilon_{\sK}$ if all queries are v0
queries, and at most the bound of Theorem~\ref{thm:v1kpre} in general.
With an untrusted commitment, $M^\ast$ of Theorem~\ref{thm:untrusted}
replaces $M$ in the definition of $y^\ast$, and the bound is that of Theorem~\ref{thm:untrusted}.
In v0, $\varepsilon$ does not depend on $G$: $\mathcal G^{(G)}$ has the weight operations,
operation inputs and commitment of $\mathcal G$. The token rule and the sampler add no term.
\end{theorem}

\begin{proof}
Suppose that the verifier accepts and that $\hat Z=M_{\mathcal G^{(G)}}(s)$. We show by induction
that $y_i=y^\ast_i$ for every $i\le G$. Let $y_1,\dots,y_{i-1}$ equal $y^\ast_1,\dots,y^\ast_{i-1}$.
By Lemma~\ref{lem:causal} with $n=G$ and $j=i-1$, column $i-1$ of $\hat Z$ is
$M_{\mathcal G}(s_0\cdots s_{P+i-2})=M_{\mathcal G}(x\,\|\,y^\ast_1\cdots y^\ast_{i-1})=\ell_i$. The
token rule for $i<G$, and the definition of $y_G$ for $i=G$, give
$y_i=R(\hat Z_{:,i-1},P+i-2)=R(\ell_i,P+i-2)=y^\ast_i$; by Lemma~\ref{lem:sampler}, $R$ is a
function, which is all that is used here. Hence an accepted response other than $y^\ast$ has
$\hat Z\neq M_{\mathcal G^{(G)}}(s)$: the verifier accepted a wrong output of $\mathcal G^{(G)}$ on
$s$. The query $s$ and the claims are fixed in the prover's first message, before any challenge,
and Theorem~\ref{thm:sound} holds for every input, so also for an $s$ that the prover chooses; the
same holds for Theorems~\ref{thm:untrusted} and~\ref{thm:v1}, and Lemma~\ref{lem:causal} holds for
$M^\ast$, whose weight products are also position-uniform. Under Fiat--Shamir the prover chooses
$s$ adaptively, which Theorem~\ref{thm:fstree} covers because the transcript absorbs the statement
first. The bound of Theorem~\ref{thm:sound} depends on $\mathcal J$, the $k_j$, $n_j$, $t_j$ and
$r$, which \texttt{with\_lm\_positions} leaves unchanged: it changes only the slices.
\end{proof}

\noindent In V1 the terms of Theorem~\ref{thm:v1} that grow with the number of leaves grow with
$T=P+G-1$, as for any prompt of that length.

\begin{corollary}[the client's seed]
\label{cor:seed}
Fix $x$, $a$ and $k$, and let the client draw the seed uniformly, independently of everything the
provider holds before the query. For every set $W$ of responses, the probability that the verifier
accepts a response in $W$ is at most $\Pr_{\mathit{seed}}[y^\ast(x,\sigma)\in W]+\varepsilon$.
\end{corollary}

\begin{proof}
The event ``accepted and in $W$'' lies in the union of $\{y^\ast\in W\}$ and of ``accepted and
different from $y^\ast$''. The reference response is a function of $M$, $x$ and $\sigma$, so the
first event has probability $\Pr_{\mathit{seed}}[y^\ast(x,\sigma)\in W]$; the second has probability
at most $\varepsilon$ by Theorem~\ref{thm:gen}.
\end{proof}

\noindent The corollary bounds the probability of every set of responses from above. It does not
prevent a provider from refusing to answer after it sees the seed, which lowers the probability of
the responses it dislikes; the distribution conditioned on acceptance can therefore differ from the
sampler's. It also says nothing about how close $R_\sigma$ is to sampling at temperature $\tau$
(Section~\ref{sec:ext-sampling}).

\para{The seed is part of the statement.} The verifier's decision depends on $\sigma$ through the
token rule, so $\sigma$ belongs to the statement. Under Fiat--Shamir, (R3) requires the hash input to
contain the whole statement, and \texttt{pvi.v0.query.absorb\_statement} absorbs
$\texttt{pvi/sampler/v1}\,\|\,a\,\|\,k\,\|\,\mathit{seed}$ (\texttt{TokenSampler.digest}), with the
graph digest, which includes the number $G$ of LM-head positions (\texttt{meta["lm\_positions"]}).
Without it, one transcript would serve the statements of every seed, the weak Fiat--Shamir
pattern~\cite{bernhard2012pitfalls,dao2023weakfs}. In mode $\Kpre$ of v0 no challenge is derived
from a transcript (Proposition~\ref{prop:kpre-fs}); the verifier holds $\sigma$ and applies the
token rule itself. V1 takes the sampler as well (\texttt{run\_cut\_query}), and its LM head is
clear, so the token rule acts on clear claims; ``an output that violates the token rule'' in
Theorem~\ref{thm:v1} is the event of Theorem~\ref{thm:gen}. Convincing a third party that the seed
was random needs a jointly generated seed, which we do not provide.

\subsection{The Fiat--Shamir Transcript: Tree Absorption and Mode \texorpdfstring{$\Kpre$}{Kpre}}
\label{app:fstree}

\para{The transcript.} For a label $L$ and a message $m$ let
$\mathrm{rec}(L,m)=\langle|L|\rangle_2\,\|\,L\,\|\,\langle|m|\rangle_8\,\|\,m$, with big-endian
lengths of 2 and 8 bytes. The transcript string starts with \texttt{pvi/fullcheck/v3}, and a
\emph{direct} absorption of $(L,m)$ appends $\mathrm{rec}(L,m)$ (\texttt{Challenger.absorb}). A
\emph{tree} absorption (\texttt{Challenger.absorb\_large}, under Fiat--Shamir, for
$|m|\ge2^{24}$ bytes) cuts $m$ into the $c=\lceil|m|/2^{22}\rceil$ chunks $m_0,\dots,m_{c-1}$ of
$2^{22}$ bytes (the last one possibly shorter), computes the leaves
\[
  \lambda_j=H_{256}\big(\texttt{pvi/leaf/v1}\,\|\,\langle|L|\rangle_2\,\|\,L\,\|\,\langle j\rangle_8
  \,\|\,m_j\big),
\]
on a thread pool, in order, and appends
$\mathrm{rec}\big(L\|\texttt{/tree},\ \langle c\rangle_8\|\langle2^{22}\rangle_8\|\langle|m|\rangle_8
\|\lambda_0\|\cdots\|\lambda_{c-1}\big)$. A message shorter than $2^{24}$ bytes is absorbed directly
under $L$. The key of a challenge with label $L_{\mathrm{ch}}$ is
$K=H_{256}(\text{transcript string}\,\|\,\texttt{challenge/}\,\|\,L_{\mathrm{ch}})$
(\texttt{Challenger.\_key}), and the challenge is expanded from $K$ as in
Appendix~\ref{app:the-conv}, with inputs $K\,\|\,\cdot$. A claims message of about \ProofLlamaLongC{}
(Llama-2-7B at 2{,}048 tokens, Section~\ref{sec:v1}) has about \TreeLeavesLlama{} leaves, so the transcript
absorbs about \TreeKBLlama{}\,kB instead of \ProofLlamaLongC{}, and the leaves are hashed in parallel.

\para{Labels.} Two conditions on labels make the transcript decodable.
(L1)~Every label, absorbed or of a challenge, is shorter than $25{,}448$ bytes, the number whose
2-byte big-endian encoding is \texttt{ch}, the first two bytes of \texttt{challenge/}.
(L2)~No label that is absorbed directly equals $L\|\texttt{/tree}$ for a label $L$ that is passed to
\texttt{absorb\_large}.
Both hold for the implementation, whose labels have a few dozen bytes at most: \texttt{absorb\_large}
receives \texttt{claims/PVC3}, \texttt{rows/I8}, \texttt{u/F31}, \texttt{cut/exc} and
\texttt{cut/mult}; the direct labels are \texttt{graph}, \texttt{params}, \texttt{op/}$\cdot$,
\texttt{col/}$\cdot$, \texttt{group/}$\cdot$, \texttt{table/}$\cdot$, \texttt{lookups}, \texttt{x},
\texttt{sampler}, \texttt{cut}, \texttt{claim/}$\cdot$, \texttt{lookup/}$\cdot$, \texttt{u/}$\cdot$
and the GKR labels \texttt{cut/}$b$\texttt{/top}, \texttt{cut/}$b$\texttt{/g}$k.j$ and
\texttt{cut/}$b$\texttt{/v}$k$ for integers $b,k,j$; the setup proof's transcript absorbs
\texttt{setup/params}, \texttt{setup/graph}, \texttt{setup/group/}$\cdot$ and \texttt{setup/msg/}$\cdot$
directly. (L2) would fail only for a weight operation named \texttt{F31/tree}.

\begin{lemma}[decoding the transcript]
\label{lem:tree-dec}
Assume (L1) and (L2). (a)~A challenge input
$w=\text{transcript string}\|\texttt{challenge/}\|L_{\mathrm{ch}}$ determines the sequence of
records and $L_{\mathrm{ch}}$. (b)~It determines, for each message, whether it was absorbed directly
or as a tree, and the message itself, or $(c,|m|,\lambda_0,\dots,\lambda_{c-1})$ respectively.
(c)~If two different messages $m\neq m'$, absorbed as trees under the same label, give the same
$(c,|m|,\lambda_0,\dots,\lambda_{c-1})$, then two different inputs of $H_{256}$ have the same
output.
\end{lemma}

\begin{proof}
(a)~After the 16-byte domain string, the next two bytes are either the length of a record's label,
which by (L1) is below $25{,}448$ and so differs from \texttt{ch}, or the bytes \texttt{ch} that start
the suffix. A record is self-delimiting through its two lengths, so reading record after record,
and stopping at \texttt{ch}, recovers the records and then $L_{\mathrm{ch}}$.
(b)~A record whose label is $L\|\texttt{/tree}$ for a label $L$ of \texttt{absorb\_large} is, by (L2),
a tree record; its payload consists of three 8-byte integers and $c$ leaves of 32 bytes. Every other
record is a direct one and contains its message.
(c)~Equal lengths give equal chunk boundaries. Since $m\neq m'$, some chunk $m_j\neq m'_j$; the two
leaf inputs share the prefix up to $\langle j\rangle_8$ and differ in the chunk, and their outputs
$\lambda_j$ are equal.
\end{proof}

\begin{theorem}[Fiat--Shamir with tree absorption]
\label{thm:fstree}
Model SHAKE-256 as a random oracle and assume (L1) and (L2). Let $\Pi$ be a public-coin protocol of
the implementation (v0 in mode $\sC$ or $\sK$, with or without the compact encoding, or the
public-coin part of V1) with $k$ challenge rounds and round-by-round error
$\varepsilon_{\mathrm{rbr}}=\max_i\varepsilon_i$, and let $\mathrm{FS}(\Pi)$ be its non-interactive
form with the transcript above, which absorbs the statement first. For every prover that makes at
most $Q$ oracle queries and outputs a statement and a proof, the probability that the statement is
false and the verifier accepts is at most
\[
  (Q+k)\,\varepsilon_{\mathrm{rbr}}+N^2\,2^{-256},
\]
where $N=Q+Q_{\mathrm{ver}}$ and $Q_{\mathrm{ver}}$ is the number of oracle queries the verifier
makes on the proof (leaves, keys and expansions).
\end{theorem}

\begin{proof}
Consider all oracle queries of the prover and of the verifier, at most $N$, in the order they are
made. Every challenge input of a query begins with the same 32 bytes $\sigma_0$: the domain string, the
framing of the first record, \texttt{graph}, and the first byte of the graph digest. (The setup proof's transcript
uses the same domain string but begins with the record \texttt{setup/params}; its challenge inputs begin with
other fixed 32 bytes $\sigma_0'$, and the argument below applies to it with $\sigma_0'$ in place of
$\sigma_0$.) Call a 32-byte value
inside a query a \emph{component} if the query begins with $\sigma_0$ and the value is a leaf in a
tree record of it, or if the query does not begin with $\sigma_0$ and the value is its first 32 bytes
(for an expansion input, its key). Call a component \emph{resolved} at the time of the query if it is
the 32-byte output of an earlier query.
Let $E_0$ be the event that some output equals $\sigma_0$, $E_1$ the event that two different inputs
have the same 32-byte output, and $E_2$ the event that, for some query, its first unresolved
component equals the output of a later query. $\Pr[E_0]\le N\,2^{-256}$. Each of the $N(N-1)/2$ pairs
of queries collides with probability $2^{-256}$, so $\Pr[E_1]\le N(N-1)/2^{257}$; each query has at
most one first unresolved component, fixed when the query is made, which a later output hits with
probability $2^{-256}$, so $\Pr[E_2]\le N(N-1)/2^{257}$. Hence $\Pr[E_0\cup E_1\cup E_2]\le
N^2\,2^{-256}$, and we assume from now on that no such event occurs.
Let $w$ be a challenge input that the verifier uses for the final proof, at the time the prover
first queried it (if the prover never did, the verifier's own query is the first). Every leaf of
$w$ is the output of the verifier's leaf query on the actual chunk. If some component of $w$ was
unresolved at that time, the first one would be the output of a later query, which is $E_2$; so all
were resolved, and by $\neg E_1$ the earlier query with that output is unique and equals the
verifier's leaf input, chunk included. With Lemma~\ref{lem:tree-dec}, $w$ therefore determined at
that time a unique partial transcript $\tau(w)$ of $\Pi$ (the statement and the messages so far)
and the label of the challenge. Conversely $\tau(w)$ and the label determine $w$, since its leaves
are the oracle outputs on its chunks. By $\neg E_0$, $K=H_{256}(w)$ differs from $\sigma_0$, so $K$
is the component of every query $K\|\cdot$; by $\neg E_1$ and $\neg E_2$, no expansion input
$K\|\cdot$ was queried before $K$ was output, so the expanded challenge is uniform and
independent of all earlier answers; the challenges of one round, drawn under distinct labels from
the same transcript string, are independent for the same reason.
As in the reduction behind (R3), the prover's queries thus induce a state-restoration prover for
$\Pi$ with at most $Q$ moves: the first query on a challenge input $w$ whose components are resolved
makes the move $\tau(w)$, whose answer fixes all challenges of that round at once (queries for the
other labels of the round make the same move), and the oracle's answers on $w$ and on the expansion
inputs are sampled jointly with that challenge; the remaining queries are answered with fresh
randomness that no accepted proof uses. The verifier's own key queries on the final proof add at
most $k$ moves. By (R2) a false statement is then accepted with probability at most
$(Q+k)\,\varepsilon_{\mathrm{rbr}}$, which with $\Pr[E_0\cup E_1\cup E_2]$ gives the bound.
\end{proof}

\begin{corollary}[v0 under Fiat--Shamir]
\label{cor:fs-v0}
For v0 in mode $\sC$, with the state function of Appendix~\ref{app:sound}, $k=2$ (the vectors
$\chi,\chi'$, then the columns) and
$\varepsilon_{\mathrm{rbr}}=\max\big(p^{-r},\max_j\prod_{i<t_j}\tfrac{k_j-1-i}{n_j-i}\big)\le
\varepsilon_{\sC}$, so a prover with $Q$ queries succeeds with probability at most
$(Q+2)\,\varepsilon_{\sC}+N^2\,2^{-256}$. For V1, $\varepsilon_{\mathrm{rbr}}$ is the largest of
these terms and of the per-round errors of Appendix~\ref{app:v1}, $(\Nleaf+|\Vtab|)/(p^8-p^2)$,
$2/p^8$ and $3/p^8$, where a round holds the challenges drawn from one transcript string: the
implementation draws $\theta_{k-1}$ and $\eta_k$ of a GKR instance together, and $\alpha$ with the
vectors $\chi$ of the clear operations.
\end{corollary}

\noindent The bound of Section~\ref{sec:fs}, $Q\varepsilon_{\mathrm{rbr}}+O(Q^22^{-256})$, should be read with
$Q+k$ in place of $Q$, or with $Q$ counting the verifier's $k$ key queries. At $Q=2^{64}$ the term
$N^2\,2^{-256}$ is about $2^{-128}$, as large as the target, because keys and leaves have 32 bytes;
a collision in the SHA-256 Merkle trees, which Theorems~\ref{thm:sound} and~\ref{thm:untrusted}
exclude, has probability of the same order (Remark~\ref{rem:rbr-max}).

\begin{remark}[relation to tree hashing]
\label{rem:tree}
The tree absorption is a one-level tree whose leaves and final node are domain-separated
(\texttt{pvi/leaf/v1} against the transcript's \texttt{pvi/fullcheck/v3}), with the leaf index, the
message's label, the chunk size, the leaf count and the length bound into it. This is the kind of
mode for which Bertoni et al.\ give sufficient conditions for soundness~\cite{bertoni2014tree};
KangarooTwelve is a deployed instance of such a mode. The proof above does not use their theorem:
it needs only Lemma~\ref{lem:tree-dec} and the birthday bound on 32-byte outputs. The index
$\langle j\rangle_8$ prevents the reordering of chunks, the label prevents the reuse of a leaf in
another message, and the length fixes the boundary of the last chunk; the chunk size is a constant
and is absorbed for versioning. \texttt{tests/test\_fs\_tree.py} checks that a flipped bit in the
first, a middle or the last chunk, one byte more or less, another label, and direct absorption of
the same bytes each change the challenge, and that the threaded leaves equal a serial reference.
\end{remark}

\begin{proposition}[mode $\Kpre$ of v0 absorbs nothing]
\label{prop:kpre-fs}
In mode $\Kpre$ of v0 the proof is one message, the claims, and no challenge is derived from a
transcript; \texttt{\_claims\_message} absorbs nothing for $\Kpre$, and the verifier uses its
precomputed $(\chi_\ell,u_\ell)$ (\texttt{Verifier.\_pre}). For every prover, with any number $Q_H$
of oracle queries, the probability that some of $Q_V$ verdicts accepts a false output is at most
$Q_V\varepsilon_{\sK}+n_{\mathrm{key}}Q_H\,2^{-256}$, where $\varepsilon_{\sK}=\sum_\ell p^{-r}$ and
$n_{\mathrm{key}}$ is the number of secret keys (one per weight operation not looked up). There is
no Fiat--Shamir term and no grinding term.
\end{proposition}

\begin{proof}
The verifier's only use of the hash function is the expansion of its own secret keys into
$\chi_\ell$ (\texttt{Verifier.precompute} with an interactive \texttt{Challenger}). Except with
probability $n_{\mathrm{key}}Q_H\,2^{-256}$ no query of the prover begins with one of these keys, and
then the prover's oracle answers are independent of the $\chi_\ell$, which are uniform. Every check
of the verifier is a function of the claims, the statement, its weights and the $\chi_\ell$, so
the prover's oracle queries give it no information on the $\chi_\ell$, and the bound of
Theorem~\ref{thm:sound} for $\Kpre$ (proved in Appendix~\ref{app:v1kpre}, with no public coins)
applies to its $Q_V$ verdicts.
\end{proof}

\noindent A $\Kpre$ proof is not transferable: only the holder of $\chi$ can check it, and the
verifier knows its own statement, so no weak Fiat--Shamir issue arises. The parameters that
\texttt{params\_for} chooses under Fiat--Shamir add 64 bits to $\beta$ although $\Kpre$ has nothing
to grind; for $\Kpre$ they act as a margin on $Q_V$ (Appendix~\ref{app:stateful}). A V1 query in
$\Kpre$ still derives $\alpha$ and the GKR challenges from the transcript
(\texttt{pvi.v1.protocol.\_query} absorbs the statement, the cut plan's digest, the first message and
the GKR messages); Corollary~\ref{cor:v1kpre-fs} bounds it.

\subsection{The Compact Encoding of the Proof}
\label{app:codec}

\para{The wire protocol.} With \texttt{wire=True} every prover message travels as bytes, which the
verifier decodes before it checks anything: the claims in the format \texttt{PVC3}
(\texttt{claimcodec.decode}, \texttt{decode\_torch}, or \texttt{decode\_device} on a GPU), the rows
of a plan's lookup tables as int8 bytes (\texttt{unpack\_rows}), $u$, the opened columns and the
folds of V1 at 31 bits per element (\texttt{unpack\_field} through \texttt{field\_message}), and in
V1 the exceptions (\texttt{pvi.v1.protocol.unpack\_exceptions}), the multiplicities (one \texttt{PVC3}
row) and each GKR transcript (\texttt{pvi.v1.gkr.Transcript.from\_bytes} with
\texttt{extfield.unpack}). For a protocol $\Pi$ of Sections~\ref{sec:defence} and~\ref{sec:v1},
which sends integer and field tensors, let $\Pi_{\mathrm{wire}}$ be the protocol in which the prover
sends bytes $b_i$ in place of message $i$ and the verifier runs the verifier of $\Pi$ on the decoded
messages, rejecting if a decoding fails. Under Fiat--Shamir the transcript of
$\Pi_{\mathrm{wire}}$ absorbs, before the next challenge, the bytes $b_i$ of every message that it
decodes (\texttt{claims/PVC3}, \texttt{rows/I8}, \texttt{u/F31}, \texttt{cut/exc},
\texttt{cut/mult}), or for the GKR rounds the decoded field elements; in each case the very buffer
that the verifier decodes (\texttt{\_claims\_message} and \texttt{decoded\_claims},
\texttt{\_fold\_message} and \texttt{open\_message}, \texttt{pvi.v1.protocol.\_query}).

\begin{lemma}[the decoders]
\label{lem:codec-dec}
Each decoder is a deterministic function $\mathrm{Dec}_i(b;c)$ of the bytes $b$ and of a public
context $c$ that is fixed before $b$ is sent: the graph, the query's shape, the parameters, the
commitment plan and the cut plan. It returns either a message of the shapes that $c$ prescribes or
a rejection. Decoded claims are integers with $|z|<2^{30}$, decoded field-packed values lie in
$[0,2^{31})$, and decoded rows lie in $[-128,127]$. Moreover the codec is designed so that
$\mathrm{Dec}_i(\mathrm{Enc}_i(m);c)=m$ for every honest message $m$ (completeness; tested, not
proved here, and not used for soundness).
\end{lemma}

\begin{proof}
The decoders read only $b$ and $c$ and draw no randomness. Their rejection rules are those listed
in the documentation of \texttt{claimcodec}: a bad magic or op count, a column count other than the
expected one (checked before anything is allocated), other flag bits, slot widths above 30, bases
outside $(-2^{30},2^{30}-2^B]$, truncated or trailing bytes, non-zero padding, Rice parameters above
32, exception positions out of order or out of range, exceptions reaching $|x|\ge2^{30}$; a wrong
length for the field-packed values and the rows; coefficients of $\Fext$ outside $[0,p)$; and, for
the exceptions frame, a bad magic, counts above the operation's entries or a wrong length. A
rejection is reported as \texttt{ClaimCodecError}, \texttt{ValueError} or \texttt{None}, which
\texttt{run\_query} and \texttt{run\_cut\_query} turn into the rejection labels
\texttt{range\_or\_shape}, \texttt{freivalds}, \texttt{columns\_shape}, \texttt{cut\_exception},
\texttt{cut\_multiplicities}, \texttt{cut\_gkr} and \texttt{cut\_final}. The value ranges follow from
these rules. That every input is either decoded or rejected, with no other exception, is the
documented contract of the decoders, which \texttt{tests/test\_claimcodec.py} checks with random
corruptions, together with the round trip and with the GPU decoder rejecting exactly what the host
decoder rejects. Soundness uses only determinism and that no query is accepted without a successful
decoding: an unexpected exception stops the verifier without acceptance.
\end{proof}

\begin{theorem}[the compact encoding preserves soundness]
\label{thm:codec}
(a)~Interactively, for every prover $P_{\mathrm{wire}}$ of $\Pi_{\mathrm{wire}}$ there is a prover $P$
of $\Pi$ such that, on every statement and for every choice of the verifier's coins,
$\Pi_{\mathrm{wire}}$ accepts if and only if $\Pi$ accepts, and every verdict of
$\Pi_{\mathrm{wire}}$ is a function of the decoded messages. Hence the bounds of Theorems~\ref{thm:sound},
\ref{thm:untrusted} and~\ref{thm:v1}, including the reuse bounds of mode $\Kpre$, hold for
$\Pi_{\mathrm{wire}}$.
(b)~$\Pi_{\mathrm{wire}}$ has the same per-round errors as $\Pi$, so Theorem~\ref{thm:fstree} gives
it the Fiat--Shamir bound of $\Pi$, where $Q$ counts every oracle query, including those spent on
other encodings of the same values. Non-canonical encodings add no term.
\end{theorem}

\begin{proof}
(a)~$P$ runs $P_{\mathrm{wire}}$ and sends $\mathrm{Dec}_i(b_i;c)$ in place of $b_i$, or, if the
decoding rejects, a message that the verifier of $\Pi$ rejects at the same point (for example a
claim matrix of the wrong shape). In $\Pi_{\mathrm{wire}}$ each challenge is drawn after $b_i$ has
been received and depends on nothing the prover sends later, and the checks of
$\Pi_{\mathrm{wire}}$ are those of $\Pi$ applied to the decoded messages. With the same coins the two
runs therefore accept together, and a failed decoding in one run is a rejection in the other.
Decoded values outside the field are rejected by
$\Pi$'s own checks (the field check of $u$ in \texttt{\_disagrees}, of the openings in
\texttt{\_columns\_shape\_ok}, of V1's folds in \texttt{pvi.v1.protocol.\_query}), so they act as no
alias. The reuse bound of $\Kpre$ (Theorem~\ref{thm:v1kpre}) needs only that every verdict is a
function of the prover's messages and the public coins until the first bad check, and verdicts of
$\Pi_{\mathrm{wire}}$ are functions of the decoded messages, hence of the bytes.
(b)~Call a partial transcript of $\Pi_{\mathrm{wire}}$ doomed if some message in it fails to decode,
or if the decoded transcript is doomed for $\Pi$. A failed decoding stays failed, and $\Pi$'s prover
messages cannot leave a doomed state, so no prover message of $\Pi_{\mathrm{wire}}$ leaves one; a
complete doomed transcript is rejected, at the failed decoding or by $\Pi$'s verifier. After bytes
$b_1,\dots,b_i$ that decode, the next challenge is uniform and independent of $b_1,\dots,b_i$, hence
of the decoded messages, so it leaves the doomed set with probability at most $\Pi$'s per-round
error. Theorem~\ref{thm:fstree} applies, since the absorbed bytes (or decoded GKR values) determine
the decoded messages. An encoding of the same values with other bytes is a different partial
transcript, and querying its challenge is one more oracle query, counted in $Q$ exactly like any
other change of the prover's message.
\end{proof}

\noindent The encoding is not canonical: the encoder chooses widths and centring, and other choices
give other bytes for the same claims (\texttt{claimcodec}, ``Soundness''). Theorem~\ref{thm:codec}
shows that soundness does not need a canonical encoding; the decoder that recomputes the encoder's
choices would only make proofs unique, which our soundness notion does not require. Soundness does need that the
verifier absorbs what it decodes: a transcript that absorbed bytes other than those decoded, or only
part of them, would let the decoded message change after the challenge.

\subsection{Two Remarks on the Round-by-Round Argument}
\label{app:rbr}

This subsection makes two points about the round-by-round argument of
Appendices~\ref{app:sound} and~\ref{app:v1}.

\begin{lemma}[the column check needs only the order of messages]
\label{lem:col-any}
Let $B\in\F^{N\times k}$ be a matrix whose rows are committed honestly as Reed--Solomon codewords
of length $n$, the rows of $\hat B$ ($B=A_j$ in the row layout, and $B$ the transpose of the stacked
weights in the transposed layout), and let $w\in\F^N$ and $u\in\F^k$ with $u\neq w^\T B$ be fixed
before the set $C$ of $t$ distinct uniform columns is drawn. Then
$\Pr_C\big[\forall c\in C:(w^\T\hat B)_c=\enc(u)_c\big]\le\prod_{i<t}\frac{k-1-i}{n-i}$. The same
bound holds for $w\in\Fext^N$ and $u\in\Fext^k$ with $u\neq w^\T B$ when the check is made on each
of the eight base-field coordinate planes.
\end{lemma}

\begin{proof}
By linearity $w^\T\hat B=\enc(w^\T B)$, and $\enc(u)$ and $\enc(w^\T B)$ are distinct codewords,
which agree on a set of at most $k-1$ positions. This set is fixed before $C$ is drawn, and a
uniform set of $t$ distinct columns lies inside it with probability at most
$\binom{k-1}{t}/\binom{n}{t}=\prod_{i<t}(k-1-i)/(n-i)$. Over $\Fext$, write
$u=\sum_{q<8}\xi^qu^{(q)}$ and $w=\sum_{q<8}\xi^qw^{(q)}$ with coordinates in $\F$; since $B$ has
entries in $\F$, $w^\T B=\sum_q\xi^qw^{(q)\T}B$, so $u^{(q)}\neq w^{(q)\T}B$ for some $q$, and the
base-field case applies to plane $q$.
\end{proof}

\para{V1's column check under Fiat--Shamir.} The lemma uses no property of $w$ other than
being fixed before the columns; in particular $w$ need not be uniform. In V1 (mode $\sC$, row
layout), $w$ is the eq-vector of the GKR's final row point, fixed when the GKR ends, $u_\ell$ is sent
in the fold message, and the columns are drawn after it: interactively by the verifier, and under
Fiat--Shamir from the transcript after it absorbed the fold message's bytes
(\texttt{open\_message} with \texttt{u\_blob}, called by \texttt{pvi.v1.protocol.\_query}). In the
round-by-round argument, the state after the folds is doomed if a check that does not read the
columns fails (a Freivalds check of a clear operation, (F3) or (F4)), or if some committed matrix
$B$ has, in some plane, a fold or combined claim
$u\neq w^\T B$ for its combination vector $w$: $\chi_\ell$ for a clear operation and the eq-vector
for a cut one in the row layout, and $\bar x_\ell$ (V1) or $[X;\mathbf 1]\chi'$ (v0) in the
transposed layout, where $u$ is the prover's $y$ (V1) or the verifier's $Z\chi'$ (v0). Only
$\chi_\ell$ is uniform. By
Lemma~\ref{lem:col-any} the column challenge leaves this state with probability at most
$\max_j\prod_{i<t_j}(k_j-1-i)/(n_j-i)$. Step~6 of the proof of Theorem~\ref{thm:v1} is this round's
analysis.

\begin{remark}[the sum and the maximum]
\label{rem:rbr-max}
In v0 the challenge of the vectors leaves a doomed state with probability at most $p^{-r}$, since
only the matrix of the first wrong operation matters, and the column challenge with probability at
most $\max_j\prod_{i<t_j}(k_j-1-i)/(n_j-i)$, since every doomed matrix must escape on its own. So
$\varepsilon_{\mathrm{rbr}}$ is the maximum of these terms, not their sum $\varepsilon_{\sC}$, and
the Fiat--Shamir bound is $(Q+k)\,\varepsilon_{\mathrm{rbr}}$ by (R2). Interactively, (R1) gives
$p^{-r}+\max_j\prod_{i<t_j}(k_j-1-i)/(n_j-i)$; the sum over $\mathcal J$ in
Theorem~\ref{thm:sound} is an upper bound, which is not needed because the first wrong operation,
and with it the matrix $j$, is fixed before any challenge.
The parameters do not use the maximum. \texttt{params\_for} sets
$\beta=\lambda+\log_2(2n_{\mathrm{op}})$, plus 64 under Fiat--Shamir (\texttt{plans.column\_bits}),
with $n_{\mathrm{op}}\ge|\mathcal J|$ the number of weight operations, and makes every term at most
$2^{-\beta}$; this meets the sum-form bounds $\varepsilon_{\sC}\le2^{-\lambda}$ and, for
$Q\le2^{64}$, $Q\varepsilon_{\sC}\le2^{-\lambda}$. With the maximum, $(Q+k)\,2^{-\beta}\le2^{-\lambda}$
needs only $\beta=\lambda+65$ for $Q+k\le2^{65}$. For Llama-2-7B ($n_{\mathrm{op}}=226$) at
$\lambda=128$ this would lower $\beta$ from $200.8$ to $193$, leave
$r=\lceil193/\log_2p\rceil=7$ unchanged, and lower each $t_j$ by about 4\%, since $t_j$ grows about
linearly with $\beta$. We keep the conservative parameters. Neither choice covers the collision term
$N^2\,2^{-256}$ of Theorem~\ref{thm:fstree}, which is about $2^{-128}$ at $Q=2^{64}$; the total
Fiat--Shamir error at that $Q$ therefore exceeds $2^{-128}$, slightly with the conservative
parameters ($2^{-128}+2^{-136.8}$) and by a factor of 2 with $\beta=\lambda+65$, and $2^{-128}$
holds for $Q\le2^{63}$ with either.
\end{remark}

\subsection{V1 in Mode \texorpdfstring{$\Kpre$}{Kpre}}
\label{app:v1kpre}

\para{Setting.} In mode $\Kpre$ the verifier holds, for every weight operation $\ell$ that is not
looked up, a secret matrix $\chi_\ell\in\F^{N_\ell\times r}$ whose columns
$\chi_{\ell,1},\dots,\chi_{\ell,r}$ are independent and uniform, drawn from a fresh key per
operation (\texttt{Verifier.precompute}), and $u_\ell=\chi_\ell^\T A_\ell$ with rows $u_{\ell,s}$.
A v0 query checks every operation by Freivalds' equation~\eqref{eq:freivalds} with
$(\chi_\ell,u_\ell)$. A V1 query checks every clear operation the same way (check~(F1)) and every
cut operation by (F2), $\chi_{\ell,s}^\T y_\ell=u_{\ell,s}\bar x_\ell$ for $s\le r$, with the same
$\chi_\ell$ and $u_\ell$ (\texttt{pvi.v1.protocol.\_kpre\_y\_ok} reads \texttt{Verifier.\_pre}), where
$y_\ell$ is the prover's fold and $\bar x_\ell=\hat X_\ell e$. To keep $\chi$ for the secret
vectors, we write $\omega_\ell$ for the eq-vector of Section~\ref{sec:v1} (its $\chi_\ell$), so that
check~(F3) uses $S_\ell=\omega_\ell^\T y_\ell$. In $\Kpre$ the verifier evaluates rule~(P4) of the
cut set from its own weights (\texttt{CutPlan.from\_graph} with \texttt{weights}).

\para{Secret checks and their errors.} Call (F1), (F2) and v0's Freivalds check the \emph{secret
checks}. The \emph{error} of a secret check $c$ of operation $\ell$ in a query is
$E_c=Z_\ell-A_\ell[\tilde X_\ell;\mathbf 1]$ for (F1) and v0, and $E_c=y_\ell-A_\ell\bar x_\ell$ for
(F2), with $\tilde X_\ell$ the derived input. It is a function of the statement, the prover's
messages and the public challenges of the query, and not of $\chi_\ell$.

\begin{lemma}[secret checks]
\label{lem:secret-check}
A secret check passes if and only if $\chi_\ell^\T E^{(q)}_c=0$ for every coordinate plane $q$ of
$E_c$ (one plane over $\F$). Hence it passes if $E_c=0$, and if $E_c\neq0$ and $E_c$ is independent
of $\chi_\ell$, it passes with probability at most $p^{-r}$.
\end{lemma}

\begin{proof}
For Freivalds' check, $\chi_\ell^\T Z_\ell-u_\ell[\tilde X_\ell;\mathbf 1]
=\chi_\ell^\T(Z_\ell-A_\ell[\tilde X_\ell;\mathbf 1])$. For (F2),
$\chi_{\ell,s}^\T y_\ell-u_{\ell,s}\bar x_\ell=\chi_{\ell,s}^\T(y_\ell-A_\ell\bar x_\ell)$ in
$\Fext$, and since $\chi_{\ell,s}$ and $A_\ell$ have entries in $\F$ this vanishes exactly when it
vanishes on every coordinate plane. If $E_c\neq0$, some plane has a nonzero column $d\in\F^{N_\ell}$,
and $\chi_{\ell,s}^\T d=0$ holds with probability $1/p$ for each of the $r$ independent uniform
columns $\chi_{\ell,s}$, so with probability $p^{-r}$ for all of them.
\end{proof}

\begin{theorem}[$\Kpre$ with reused secret vectors, v0 and V1]
\label{thm:v1kpre}
Let $Q$ queries use the same secret vectors: any mix of v0 and V1 queries, generation queries
included, chosen adaptively, with the verdict of each one revealed before the next. The probability
that some query accepts a false statement is at most
\[
  \sum_{i\le Q}(\varepsilon^{\mathrm{lu}}_i+\varepsilon^{\mathrm{gkr}}_i)\;+\;Q\sum_\ell p^{-r},
\]
where $\varepsilon^{\mathrm{lu}}_i=(\Nleaf+|\Vtab|)/(p^8-p^2)$ and $\varepsilon^{\mathrm{gkr}}_i=\sum_\gamma\varepsilon_{\mathrm{GKR}}(n_\gamma)$
for a V1 query, with the values of its shape, $\varepsilon^{\mathrm{lu}}_i=\varepsilon^{\mathrm{gkr}}_i=0$ for a v0 query, and $\ell$ ranges
over the weight operations that are not looked up. For queries of one shape this is $Q$ times the
$\Kpre$ bound of Theorem~\ref{thm:v1}, and for v0 alone it is $Q\varepsilon_{\sK}$, the reuse bound
of Theorem~\ref{thm:sound}.
\end{theorem}

\begin{proof}
\emph{The hybrid.} Let $\mathsf H$ be the experiment in which the verifier replaces the outcome of
every secret check $c$ by the predicate $[E_c=0]$ and otherwise runs unchanged; it reports verdicts
in the same way. Call a secret check \emph{bad} if it is evaluated, $E_c\neq0$, and the real check
would pass. Run the real experiment and $\mathsf H$ with the same coins of the prover, the same public
challenges and the same $\chi$. By Lemma~\ref{lem:secret-check}, a secret check that is not bad has
the same outcome in both: it passes in both if $E_c=0$, and fails in both otherwise. Since the
verifier checks in a fixed order and stops at the first failure, the two runs are identical, verdicts
and labels included, up to the first bad check, and a bad check occurs in one run exactly when it
occurs in the other.

\emph{Bad checks.} In $\mathsf H$ the verifier never uses $\chi$, so the prover's whole view, every
message and every public challenge is independent of $\chi$. For query $i$ and operation $\ell$
there is at most one secret check, and its error $E_{i,\ell}$ is a function of that view;
conditioned on the view, Lemma~\ref{lem:secret-check} bounds the probability that it is bad by
$p^{-r}$. A union bound over $i\le Q$ and $\ell$ gives $\Pr_{\mathsf H}[\text{bad}]\le
Q\sum_\ell p^{-r}$.

\emph{False acceptance without a bad check.} In $\mathsf H$ each query is checked by a public-coin
verifier $V_{\mathrm{pub}}$, whose secret checks are the deterministic predicates $[E_c=0]$. We show
that, conditioned on everything before query $i$, $V_{\mathrm{pub}}$ accepts a false statement in
query $i$ with probability at most $\varepsilon^{\mathrm{lu}}_i+\varepsilon^{\mathrm{gkr}}_i$. Follow the proof of Theorem~\ref{thm:v1}
(Appendix~\ref{app:v1}). By its Step~1 there is a first locally incorrect operation $\ell^\ast$,
whose derived input is honest. If $\ell^\ast$ is clear (always, in v0), its claim is wrong and, by
the range checks, $E\neq0$, so $[E=0]$ fails. If $\ell^\ast$ is cut, Lemma~\ref{lem:v1-nowrap} gives
a leaf outside $\Vtab$, which (with the multiplicities) is fixed before $\alpha$; by
Lemma~\ref{lem:v1-logup} and the pole count of Step~4, the logUp identity holds at the fresh
$\alpha$ with probability at most $\varepsilon^{\mathrm{lu}}_i$; otherwise some instance has a wrong root, which the
GKR carries to wrong final claims except with probability at most $\varepsilon^{\mathrm{gkr}}_i$
(Lemma~\ref{lem:v1-gkr}). Check~(F3) compares the final claims with values computed from
$S_\ell=\omega_\ell^\T y_\ell$; if $y_\ell=A_\ell\bar x_\ell$ for every cut operation, these are the
true values (Lemma~\ref{lem:v1-final}) and (F3) fails; otherwise some $[E_{(F2)}=0]$ fails. The
public challenges of query $i$ are fresh, so these bounds hold conditioned on the history.

\emph{Conclusion.} The probability that the real experiment accepts some false statement is at most
the probability of a bad check plus the probability of a false acceptance without one. Both are the
same in $\mathsf H$, by the coupling, and the two bounds above give the theorem. Throughout, a verdict
of $\mathsf H$, and of the real run before the first bad check, is a function of the prover's
messages and the public challenges, so revealing it, and the label of the failed check, gives the
prover nothing about $\chi$.
\end{proof}

\begin{corollary}[V1 in $\Kpre$ under Fiat--Shamir]
\label{cor:v1kpre-fs}
Under Fiat--Shamir, let the prover make $Q_H$ oracle queries and submit $Q_V$ proofs, each with at
most $k$ challenge rounds. Some false statement is accepted with probability at most
\[
  (Q_H+kQ_V)\,\varepsilon_{\mathrm{pub}}+Q_V\sum_\ell p^{-r}+N^2\,2^{-256}+n_{\mathrm{key}}Q_H\,2^{-256},
\]
where $\varepsilon_{\mathrm{pub}}=\max\big((\Nleaf+|\Vtab|)/(p^8-p^2),3/p^8\big)$ is the largest
per-round error of the public rounds, $N$ is $Q_H$ plus the number of the verifier's oracle queries
on all submitted proofs, and $n_{\mathrm{key}}$ is as in Proposition~\ref{prop:kpre-fs}.
\end{corollary}

\begin{proof}
The hybrid and the bound on bad checks are as above: the oracle is independent of $\chi$ except with
probability $n_{\mathrm{key}}Q_H2^{-256}$ (Proposition~\ref{prop:kpre-fs}). In $\mathsf H$ the prover
interacts with $\mathrm{FS}(V_{\mathrm{pub}})$ only. $V_{\mathrm{pub}}$ has the public rounds of V1,
whose per-round errors are those of the Fiat--Shamir paragraph of Appendix~\ref{app:v1}, and checks
the deterministic predicates $[E_c=0]$ in place of the secret checks; a state function for it is that
of the proof of Theorem~\ref{thm:v1}, with the predicates in place of (F1) and (F2). The proof of
Theorem~\ref{thm:fstree}, applied with $Q_V$ final proofs, gives a state-restoration prover with at
most $Q_H$ moves plus $k$ moves per submitted proof. (R2) is stated for one final proof; its argument
is a union bound over the moves, each of which leaves a doomed state with probability at most
$\varepsilon_{\mathrm{pub}}$, and an accepted proof of a false statement needs one such move, so the
same argument bounds the success with $Q_V$ final proofs by $(Q_H+kQ_V)\,\varepsilon_{\mathrm{pub}}$.
\end{proof}

\subsection{Reusing the Secret Vectors: the Stateful Verifier of Mode \texorpdfstring{$\Kpre$}{Kpre}}
\label{app:stateful}

The $\Kpre$ verifier keeps state across queries: the secret pair $(\chi_\ell,u_\ell)$ of every
operation and, for the bound of Theorem~\ref{thm:v1kpre} to be usable, the number of verdicts it has
given with it. The bound grows linearly in that number, and the factor cannot be removed.

\begin{proposition}[the factor $Q$ is necessary]
\label{prop:qtight}
Consider next-token queries (one LM-head position, so no token rule), and let the weight operation
$\ell$ whose claims are the graph's output (for a decoder, the LM head) have $N_\ell\ge Q$ rows and be
checked in $\Kpre$ with $r$ secret columns. There is a prover that, without looking at the verdicts,
makes $Q$ queries with false outputs such that at least one is accepted with probability
$1-(1-p^{-r})^Q\ge Qp^{-r}-\binom Q2p^{-2r}$. After its first accepted query, it can make every later
next-token query accept a false output with probability~1.
\end{proposition}

\begin{proof}
Fix distinct rows $\rho_1,\dots,\rho_Q$ of operation $\ell$. In query $i$, on any input, the prover
computes honestly and changes the claim $z$ of $\ell$ in row $\rho_i$ (and the query's one column) by
$\delta_i=-\mathrm{sign}(z)$, or by $+1$ if $z=0$, which keeps it inside the range check. Since $\ell$
is the last weight operation, every derived input stays honest, every other check passes, and the
output is wrong in row $\rho_i$. The error of $\ell$'s check is $\delta_i$ times the matrix with a
single one in row $\rho_i$, so the check passes if and only if row $\rho_i$ of $\chi_\ell$ is zero,
which has probability $p^{-r}$. Distinct rows of $\chi_\ell$ are independent, so the $Q$ events are
independent, which gives the first claim; the inequality is Bonferroni's. If query $i$ is accepted,
row $\rho_i$ of $\chi_\ell$ is zero, and every later change of the claims in row $\rho_i$ alone, of
any size within the range, passes the check and makes the output false.
\end{proof}

\begin{remark}[using a $\Kpre$ key]
\label{rem:stateful}
(a)~\emph{Key lifetime.} The verifier must count the verdicts of each key and replace the key (a new
precomputation) before $Q_V\varepsilon_{\sK}$ exceeds its target, or choose
$r=\lceil(\lambda+\log_2Q_{\max}+\log_2n_{\mathrm{op}})/\log_2p\rceil$ for a fixed maximum
$Q_{\max}$. For Llama-2-7B, with its 226 weight operations, the interactive parameters at
$\lambda=128$ ($r=5$) give $\varepsilon_{\sK}=226\,p^{-5}=2^{-146.7}$, which keeps
$Q_V\varepsilon_{\sK}\le2^{-128}$ for $Q_V\le2^{18.7}$ verdicts only, about $4\cdot10^5$. The
parameters with the 64-bit margin ($r=7$) give $2^{-208.5}$ and $Q_V\le2^{80.5}$.
(b)~\emph{Every verdict counts.} $Q_V$ counts the verdicts of every query that uses the key, accepted
or rejected, v0 or V1, a generated response counting as one, and summed over all verifiers that
share the key.
(c)~\emph{What the prover may observe.} The hybrid of Theorem~\ref{thm:v1kpre} covers the verdict and
the label of the failed check, because both are functions of the prover's messages until the first
bad check. It does not cover anything else that depends on $\chi$, such as a running time that
depends on the compared values. The CPU verifier does not run in constant time: it compares the two
sides of Freivalds' check operation by operation with \texttt{torch.equal} and stops at the first
failing operation (\texttt{pvi.v0.verifier.\_disagrees}), and check~(F2) returns at the first mismatch
(\texttt{pvi.v1.protocol.\_kpre\_y\_ok}). Which operation fails first is a function of the prover's
messages until the first bad check, so the hybrid covers it; the time of one comparison can depend on
where its two sides first differ, and so on $\chi$, which the proof does not cover. A verifier whose
timing the provider can observe should complete every comparison before it answers, or answer only
accept or reject after a fixed-time check.
(d)~\emph{The key must stay secret.} The bound gives nothing once $\chi_\ell$ leaks. For instance,
when $N_\ell>r\log_2p$, two different subsets of rows have the same sum of rows of $\chi_\ell$
(pigeonhole over the $2^{N_\ell}$ subsets and $p^r$ values), so there is a nonzero
$d\in\{-1,0,1\}^{N_\ell}$ with $\chi_\ell^\T d=0$. For the LM head ($N_\ell=32{,}000$ rows for
Llama-2-7B, far above $7\log_2p\approx216$), adding $d$ to a column of its claims, where this keeps
them inside the range check, passes its check and makes the output false. We have not studied how
fast such a $d$ can be found from $\chi_\ell$.
(e)~\emph{After a false acceptance.} The bound concerns the first false acceptance. As
Proposition~\ref{prop:qtight} shows, an accepted forgery can reveal enough of $\chi$ to make later
forgeries certain, and the verifier cannot detect this. A key under which a forgery may have been
accepted gives no guarantee for its later queries.
(f)~\emph{Designated verifier.} Proofs of $\Kpre$ convince only the holder of $\chi$.
\end{remark}

\subsection{Acceptance of Edits of \texorpdfstring{$k$}{k} Neurons by the Path Protocol}
\label{app:paths}

The path protocol (Section~\ref{sec:path}) catches an edit only if a path checks one of the sites it
made inconsistent. This gives an acceptance bound that depends only on the number of such sites,
not on what the edit does to the output.

\begin{proposition}[acceptance of a $k$-site edit]
\label{prop:paths}
Let a forged trace have a set $B$ of at most $k$ sites whose local check fails, all in one layer of
width $N$, every other site passing its check, and let the verifier run $s$ independent paths, each of
which checks each site of that layer with probability at most $1/N$ (a path checks the sites it visits;
the values it opens to check them are not checked). Then the verifier accepts with
probability at least $(1-k/N)^s$. Rejection with probability at least $1-2^{-\lambda}$ therefore needs
\[
  s\;\ge\;s_\lambda(k,N)=\Big\lceil\frac{\lambda\ln2}{-\ln(1-k/N)}\Big\rceil,
\]
\[
  \Big(\frac Nk-1\Big)\lambda\ln2\;\le\;s_\lambda(k,N)\;\le\;\Big\lceil\frac Nk\,\lambda\ln2\Big\rceil.
\]
RandPathTest~\cite{anchuri2026verifiable} with uniform parents meets the hypothesis for a layer whose
consumer is fully connected, such as the hidden layer of an MLP or FFN: a path that reaches the
consumer moves to a uniform site of the layer, and a path visits a layer at most once. If every path
reaches the consumer and $|B|=k$, the acceptance probability is exactly $(1-k/N)^s$.
\end{proposition}

\begin{proof}
The prover opens the values of its committed trace, so a path that checks no site of $B$ passes all
its checks, and a path that checks a site of $B$ fails there. Path $i$ therefore rejects with
probability at most $\sum_{b\in B}\Pr[\text{path $i$ checks $b$}]\le k/N$, and since the paths use
independent coins the verifier accepts with probability
$\prod_{i\le s}(1-\Pr[\text{path $i$ rejects}])\ge(1-k/N)^s$. If the rejection probability is at least
$1-2^{-\lambda}$, then $(1-k/N)^s\le2^{-\lambda}$, which gives $s\ge s_\lambda(k,N)$. The bounds on
$s_\lambda$ follow from $x\le-\ln(1-x)\le x/(1-x)$ for $0\le x<1$. In a fully connected consumer each
neuron's parents are all $N$ sites, so a uniform parent is uniform on the layer whatever the path's
history; a path moves towards the input at every step, so it visits a layer at most once, and when
every path reaches the consumer each checks one uniform site, which gives equality.
\end{proof}

\noindent For $k=1$ and $\lambda=40$ this is the $27.7N$ of Section~\ref{sec:dichotomy}, and it
reproduces the counts of Section~\ref{sec:break}: \PathsLlama{} paths for Llama-2-7B and \AttackOptLPaths{} for
OPT-6.7B. The bound holds for uniform paths. A sampler that weights parents by the claimed values
behaves differently, better against edits that raise values and not at all against edits to zero
(Appendix~\ref{ext:attack-inrange}); Proposition~\ref{prop:blind} and
Theorem~\ref{thm:dich}(a) treat samplers in general.

\begin{table}[t]
  \caption{Uniform paths needed to catch an edit of $k$ neurons in one FFN layer of width $N$ with
  probability $1-2^{-40}$: $s_{40}(k,N)$ of Proposition~\ref{prop:paths}. Computed from the formula;
  $k=1$ is the edit measured on OPT-6.7B, and $k=2$, $6$ and $60$ are the median in-range edit, the
  median lowering-only edit and the largest in-range edit (moves in either direction) measured on
  OPT-1.3B (Section~\ref{sec:break}); $k=16$, and every $k$ for the other models, is illustrative.}
  \label{tab:paths}
  \centering\footnotesize
  \setlength{\tabcolsep}{3.5pt}
  \begin{tabular}{@{}lrrrrrr@{}}
    \toprule
    Model & $N$ & $k=1$ & $k=2$ & $k=6$ & $k=16$ & $k=60$ \\
    \midrule
    OPT-1.3B    &  8{,}192 & 227{,}117 & 113{,}552 & 37{,}842 & 14{,}182 & 3{,}772 \\
    Qwen3-4B    &  9{,}728 & 269{,}704 & 134{,}845 & 44{,}940 & 16{,}844 & 4{,}482 \\
    Llama-2-7B  & 11{,}008 & 305{,}193 & 152{,}590 & 50{,}854 & 19{,}062 & 5{,}073 \\
    Llama-2-13B & 13{,}824 & 383{,}269 & 191{,}628 & 63{,}867 & 23{,}942 & 6{,}375 \\
    OPT-6.7B    & 16{,}384 & 454{,}248 & 227{,}117 & 75{,}697 & 28{,}378 & 7{,}558 \\
    \bottomrule
  \end{tabular}
\end{table}

\begin{table*}[t]
  \caption{Probability $(1-k/N)^s$ that $s$ uniform paths accept an edit of $k$ neurons in one FFN
  layer (computed from Proposition~\ref{prop:paths}).}
  \label{tab:paths-acc}
  \centering\footnotesize
  \setlength{\tabcolsep}{3.5pt}
  \begin{tabular}{@{}llrrrrr@{}}
    \toprule
    Model & $s$ & $k=1$ & $k=2$ & $k=6$ & $k=16$ & $k=60$ \\
    \midrule
    Llama-2-7B & $10^3$ & 0.913 & 0.834 & 0.580 & 0.234 & $4.2\cdot10^{-3}$ \\
               & $10^4$ & 0.403 & 0.163 & $4.3\cdot10^{-3}$ & $4.8\cdot10^{-7}$ & $1.8\cdot10^{-24}$ \\
               & $10^5$ & $1.1\cdot10^{-4}$ & $1.3\cdot10^{-8}$ & $2.1\cdot10^{-24}$ & $6.8\cdot10^{-64}$ & $<10^{-200}$ \\
    \midrule
    OPT-6.7B   & $10^3$ & 0.941 & 0.885 & 0.693 & 0.376 & 0.026 \\
               & $10^4$ & 0.543 & 0.295 & 0.026 & $5.7\cdot10^{-5}$ & $1.2\cdot10^{-16}$ \\
               & $10^5$ & $2.2\cdot10^{-3}$ & $5.0\cdot10^{-6}$ & $1.2\cdot10^{-16}$ & $3.7\cdot10^{-43}$ & $<10^{-150}$ \\
    \bottomrule
  \end{tabular}
\end{table*}
\endgroup
